\appendix

\section{Complementos sobre el acondicionamiento}
\label{anx:acond}

La figura~\ref{fig:acond_crudo} muestra los registros crudos completos, y la figura~\ref{fig:acond_espectro_crudo}, sus espectros. La deriva aparece como potencia por debajo de \SI{0.5}{\hertz}, unas 75 veces mayor post ejercicio que en reposo; el pasaaltos del módulo la atenúa, pero no la elimina. En \SI{50}{\hertz} no hay un pico por encima de las frecuencias vecinas. La figura~\ref{fig:acond_espectro_antes_despues} compara los espectros antes y después del acondicionamiento. El firmware no lee las salidas de electrodo suelto (LO+ y LO$-$), así que un electrodo suelto solo se reconocería como un tramo plano.

\begin{figure}[htbp]
  \centering
  \includegraphics[width=\textwidth]{acond_crudo_completo}
  \caption{Registros crudos en reposo y post ejercicio. La línea punteada marca \SI{1.65}{\volt}.}
  \label{fig:acond_crudo}
\end{figure}

\begin{figure}[htbp]
  \centering
  \includegraphics[width=\textwidth]{acond_espectro_crudo}
  \caption{Densidad espectral de los registros crudos (Welch, ventanas de \SI{8}{\second}). Izquierda: espectro completo; la línea punteada marca \SI{50}{\hertz}. Derecha: de 0 a \SI{3}{\hertz}; la línea punteada marca \SI{0.5}{\hertz}.}
  \label{fig:acond_espectro_crudo}
\end{figure}

\begin{figure}[htbp]
  \centering
  \includegraphics[width=\textwidth]{acond_espectro_antes_despues}
  \caption{Densidad espectral de 0 a \SI{5}{\hertz} de los registros crudos (trazo claro) y acondicionados (trazo oscuro). La línea punteada marca el corte del pasaaltos.}
  \label{fig:acond_espectro_antes_despues}
\end{figure}

\section{Complementos sobre el tacograma y el periodograma}
\label{anx:lomb}

\paragraph{Índices del tacograma.}
La pendiente de la tabla~\ref{tab:tacograma} es la de la recta ajustada por cuadrados mínimos a la frecuencia instantánea. El criterio de atípicos del \SI{30}{\percent} detecta latidos perdidos, que duplican el RR, y detecciones falsas, que lo parten en dos, pero sigue los cambios lentos de la recuperación. Con \SI{20}{\percent} se marcaban como atípicos dos intervalos reales del reposo, en las crestas de una oscilación lenta. En el registro post ejercicio, revisar solo los extremos de cada intervalo dejaba pasar dos intervalos que contenían un latido perdido.

\paragraph{Periodograma de Lomb.}
Para una serie $x_n$ de media nula, tomada en los instantes $t_n$, la potencia en $\omega = 2\pi f$ es
\begin{equation}
  P(\omega) = \frac{1}{2}\left[
    \frac{\left(\sum_n x_n \cos\omega(t_n-\tau)\right)^2}{\sum_n \cos^2\omega(t_n-\tau)} +
    \frac{\left(\sum_n x_n \sin\omega(t_n-\tau)\right)^2}{\sum_n \sin^2\omega(t_n-\tau)}
  \right],
  \qquad
  \tan(2\omega\tau) = \frac{\sum_n \sin 2\omega t_n}{\sum_n \cos 2\omega t_n},
  \label{eq:lomb}
\end{equation}
donde $\tau$ hace que los ajustes del seno y del coseno sean independientes. Si las muestras son equiespaciadas, $P$ coincide con el periodograma clásico.

\paragraph{Grilla y normalización.}
La resolución es $1/T$; la grilla cuatro veces más fina ubica mejor los picos, pero no agrega resolución. El límite de \SI{0.5}{\hertz} cubre HF con margen y queda por debajo de la mitad de la frecuencia media de latidos en todos los registros (más de \SI{0.7}{\hertz} a \SI{84}{lpm}). La densidad unilateral es $p(f) = 2P/\bar f_s$, con $\bar f_s = N/T$, de modo que el área bajo $p$ es la varianza del tacograma (teorema de Parseval). El polinomio de grado 2 solo sigue variaciones más lentas que la banda LF. Cada ventana de \SI{1}{\minute} tiene su propio tratamiento previo y su propia grilla.

\begin{table}[htbp]
  \centering
  \caption{Índices espectrales. $N$: RR usados; $T$: duración; $f_\mathrm{HF}$: pico en HF. Las tres últimas filas son ventanas de \SI{1}{\minute} del registro post ejercicio (la última dura \SI{51}{\second}).}
  \label{tab:lomb}
  \small
  \setlength{\tabcolsep}{4pt}
  \begin{tabular}{@{}lrrrrrrrr@{}}
    \toprule
    Registro & $N$ & $T$ (s) & FC (lpm) & LF (\si{\milli\second\squared}) & HF (\si{\milli\second\squared}) & LF/HF & $\mathrm{LF_{nu}}$ (\%) & $f_\mathrm{HF}$ (Hz) \\
    \midrule
    Reposo                 &  78 &  55 &  84 & 4870 & 1029 & \num{4.7} & 83 & \num{0.18} \\
    Post ejercicio, completo & 359 & 171 & 131 &   48 &   21 & \num{2.2} & 69 & \num{0.19} \\
    \quad minuto 1         & 140 &  60 & 141 &   56 & \num{4.2} & \num{13.3} & 93 & \num{0.36} \\
    \quad minuto 2         & 126 &  60 & 128 &   44 &   47 & \num{0.9} & 48 & \num{0.19} \\
    \quad minuto 3         &  93 &  51 & 121 &   23 &   16 & \num{1.5} & 59 & \num{0.21} \\
    \bottomrule
  \end{tabular}
\end{table}
